\documentclass[12pt, a4paper]{article}
\usepackage[margin=2.5cm]{geometry}
\usepackage{booktabs}
\usepackage{array}
\usepackage{tabularx}
\usepackage{xcolor}
\usepackage{colortbl}
\usepackage{enumitem}
\usepackage{titlesec}
\usepackage{parskip}
\usepackage{hyperref}
\usepackage{multirow}

\definecolor{headerblue}{RGB}{30, 90, 160}
\definecolor{rowgray}{RGB}{240, 240, 240}
\definecolor{accentblue}{RGB}{0, 70, 140}

\hypersetup{
    colorlinks=true,
    linkcolor=accentblue,
    urlcolor=accentblue
}

\titleformat{\section}{\large\bfseries\color{accentblue}}{}{0em}{}[\titlerule]
\titleformat{\subsection}{\normalsize\bfseries}{}{0em}{}

\renewcommand{\arraystretch}{1.5}

\title{%
  {\Large\bfseries Machine Learning Lab Projects}\\[0.4em]
  {\normalsize Select One Project Per Group}
}
\author{}
\date{}

\begin{document}

\maketitle
\thispagestyle{empty}

% ── Overview ────────────────────────────────────────────────────────────────────
\section{Overview}

Each group must select one project from the list below. You are expected to independently find a suitable
dataset, apply appropriate Machine Learning techniques, and demonstrate a working system at the end of
the semester.

% ── Available Projects ──────────────────────────────────────────────────────────
\section{Available Projects}

\begin{table}[ht]
\centering
\renewcommand{\arraystretch}{1.7}
\begin{tabularx}{\textwidth}{
  >{\centering\bfseries}p{0.6cm}
  >{\bfseries}p{4.5cm}
  >{\raggedright\arraybackslash}X
}
\toprule
\rowcolor{headerblue}
\textcolor{white}{\#} &
\textcolor{white}{Project Title} &
\textcolor{white}{Brief Description} \\
\midrule

1 & Network Intrusion Detection System
  & Build a smart system that identifies malicious or suspicious network traffic. Explore both supervised
    and unsupervised ML techniques, and compare their effectiveness in detecting intrusions. \\
\rowcolor{rowgray}
2 & Student Performance Prediction
  & Predict students' final grades using historical academic records, attendance, and personal background
    data. Apply regression or classification models and analyze which features influence performance
    the most. \\

3 & Student Dropout Risk Predictor
  & Develop an early-warning system that flags students likely to drop out before it is too late. Use
    academic, attendance, and demographic data and ensure your model explains why a student is at risk. \\
\rowcolor{rowgray}
4 & CI/CD Build Failure Predictor
  & Predict whether a software build will fail before it runs, using commit metadata. Expose your model
    as a REST API that can serve as a pre-build gate in real CI pipelines. \\

5 & Water Potability Classifier
  & Classify whether a water sample is safe to drink from chemical sensor readings. Identify the most
    dangerous chemical factors and deploy the model as a public API for real-world use. \\
\rowcolor{rowgray}
6 & Hospital Appointment No-Show Predictor
  & Predict which patients are likely to miss their appointment using demographics, booking info, and
    reminder data. Identify the strongest behavioral predictors to help design better intervention
    strategies. \\

7 & Industrial Predictive Maintenance
  & Predict potential mechanical failures in industrial machinery before they occur using time-series
    sensor logs. Feature engineering and high recall are key --- a missed failure is costly. \\
\rowcolor{rowgray}
8 & Real Estate / Dynamic Pricing Engine
  & Build a regression model to predict property prices from structural, location, and economic features.
    Handle skewed distributions, encode categorical variables, and evaluate using RMSE and MAE. \\

9 & Phishing / Spam Detection System
  & Develop a binary classifier to detect malicious emails or URLs using text and URL features. Apply NLP
    preprocessing, TF-IDF vectorization, and evaluate with F1-Score. \\
\rowcolor{rowgray}
10 & Collaborative Recommender System
   & Build a movie or book recommendation engine that suggests items based on user ratings. Use matrix
     factorization or KNN to handle sparse user-item interaction data. \\
\bottomrule
\end{tabularx}
\end{table}

% ── Required Steps ──────────────────────────────────────────────────────────────
\section{Required Steps (All Projects)}

Every group must complete \textbf{all} of the following, regardless of which project they choose:

\begin{itemize}[leftmargin=1.5em, itemsep=0.3em]
    \item Find a dataset.
    \item Apply ML techniques and train a model.
    \item Find the best possible model with hyperparameter tuning.
    \item Demo the project in the last lab using FastAPI or a web server.
\end{itemize}

% ── Demo Schedule ───────────────────────────────────────────────────────────────
\section{Demo Schedule}

\begin{itemize}[leftmargin=1.5em, itemsep=0.3em]
    \item Live demos will be held in the \textbf{second last lab session} of the semester.
    \item Each group will be allotted a fixed time slot for demonstration and viva.
    \item Submit \textbf{one printed report per group} on the day of your demo.
    \item \textbf{No online report submissions} will be accepted.
\end{itemize}

% ── Important Policies ──────────────────────────────────────────────────────────
\section{Important Policies}

\begin{itemize}[leftmargin=1.5em, itemsep=0.3em]
    \item \textbf{No extensions} will be granted under any circumstances.
    \item \textbf{Plagiarism or code copying} will result in zero marks for the entire group.
\end{itemize}

% ── Marks Distribution ──────────────────────────────────────────────────────────
\section{Evaluation \& Submission Criteria}

\subsection*{Marks Distribution}

\begin{table}[ht]
\centering
\begin{tabular}{lc}
\toprule
\rowcolor{headerblue}
\textcolor{white}{\textbf{Component}} & \textcolor{white}{\textbf{Marks}} \\
\midrule
Live Demo + Viva  & 20 \\
\rowcolor{rowgray}
Printed Report    & 10 \\
\midrule
\textbf{Total}    & \textbf{30} \\
\bottomrule
\end{tabular}
\end{table}

\end{document}
